% !TEX root = ../main.tex
\section{Discussion and limitations}
\label{sec:discussion}

The completed episodes demonstrate a connected execution path: live model
responses reach public deliberation and binding votes, the local contract
settles both truth cases, and local exports and centralized tracking retain
their measurements. The preceding provider-quota failure also illustrates
why infrastructure failures must remain visible. Successful execution in the
retry does not erase the two incomplete episodes in the earlier attempt.

The sample is too small for a stable estimate of population accuracy, a treatment
effect, or an absence of collusion. There is one short block, one model family,
and two clear fictional scenarios. Shared prompts, evidence, and model behavior
can explain unanimous votes. An observed SD of zero across two episodes says
only that the displayed episode values coincide; it does not establish zero
uncertainty. The seven jurors and their 21 pairs per episode are dependent
measurements within those episodes.

Several baseline controls also restrict the questions the data can answer.
Private-channel abuse is unobservable without a private opportunity, and bribe
acceptance is inapplicable without transfers. Forced claim and challenge actions
do not reveal voluntary party agreement. Cross-episode consistency under
disabled memory does not measure learned coordination. The baseline has no
auditor, so auditor--juror collusion is outside its scope.

The probes in Section~\ref{sec:probes} carry stronger caveats. Each setting
has one episode, one model and one false claim, and model sampling is
stochastic, so a repetition could differ. Each compared pair changes one
mechanism, but runs at a different time and code revision. The
evidence-consistent side was always the challenger, so no juror faced an offer
that paid for the truthful vote against an offer that paid more for a false
one. Accepting an aligned payment is therefore not evidence that money changed
a vote. In this design, the auditor's finding has no consequences, which gives
no participant a reason to buy it. The prohibition changed what agents were
told, not their payoffs, and nothing enforced it. The reasoning monitor is an
uncalibrated single-model judge whose rogue label treats accepting a permitted
payment as a breach of a juror's duty. Repeated episodes, other models, true
claims and validated monitor labels are needed before any of these
observations becomes a finding about rates or effects.

The next stage is the planned two-block, 20-episode baseline, followed by six
single-factor screens and the 64-cell design~\cite{project-design}. Comparing
evidence levels requires care: the current clear and ambiguous libraries use
different claims, so that contrast is confounded by scenario set. Additional
fresh blocks, matched scenarios, and a prespecified analysis are needed before
interpreting factor effects or interactions.
